\subsection*{\textbf{c. Expressiveness and Flexibility}}

A named dimension can appear in several shape expressions. For example, a buffer may use half of M along one axis and a padded N along another. These expressions show how the new shape depends on the input dimensions.

An anonymous dimension can still be read through the input's shape. A symbolic name is useful when that same value occurs in several expressions: changing one derivation then leaves its relationship to the other dimensions visible.

For example, you can easily express relationships between dimensions using symbolic variables:

\begin{lstlisting}[language=croqtile,caption={c. Expressiveness and Flexibility}]
mdspan sp : [M, N]; // Symbolic dimensions
mdspan spn : sp [1, M / 2, N / 4]; // Derived shape based on symbolic expressions
\end{lstlisting}

The benefit is that the dependency is written where the shape is declared. The program still has to satisfy the resulting size and divisibility requirements for each call.
